\documentclass[a4paper,11pt]{article}
\usepackage[left=59pt,right=59pt,top=51pt,bottom=54pt,footskip=31pt]{geometry}
\usepackage{amsmath,amsthm}
\usepackage{newtxtext,newtxmath}
\usepackage[hidelinks,bookmarksnumbered]{hyperref}
\hypersetup{pdftitle={Quantitative stationary crystallization for the E8 and Leech lattices},pdfauthor={Kenta Kitamura},pdfsubject={Stationary Riesz energy, quantitative lattice stability, and zero-excess laws}}
\setlength{\parindent}{0pt}
\setlength{\parskip}{7pt}
\setlength{\emergencystretch}{2em}
\AtBeginDocument{\fontsize{11.4}{16}\selectfont}
\makeatletter
\renewcommand\section{\@startsection{section}{1}{0pt}{18pt}{9pt}{\normalfont\fontsize{14.5}{18}\bfseries}}
\renewcommand\subsection{\@startsection{subsection}{2}{0pt}{11pt}{6pt}{\normalfont\fontsize{12}{16}\bfseries}}
\renewcommand\@seccntformat[1]{\csname the#1\endcsname.\quad}
\makeatother
\newtheorem{theorem}{Theorem}[section]
\newtheorem{lemma}[theorem]{Lemma}
\newtheorem{proposition}[theorem]{Proposition}
\newtheorem{corollary}[theorem]{Corollary}
\theoremstyle{remark}
\newtheorem{remark}[theorem]{Remark}
\numberwithin{equation}{section}
\newcommand{\R}{\mathbb R}
\newcommand{\Z}{\mathbb Z}
\newcommand{\E}{\mathbb E}
\newcommand{\Prob}{\mathbb P}
\newcommand{\calE}{\mathcal E}
\newcommand{\calP}{\mathcal P}
\newcommand{\vol}{\operatorname{vol}}
\newcommand{\dist}{\operatorname{dist}}
\newcommand{\Aut}{\operatorname{Aut}}
\newcommand{\Law}{\operatorname{Law}}
\newcommand{\Haar}{\operatorname{Haar}}
\newcommand{\BR}{\overline B_R}
\newcommand{\ind}{\mathbf 1}
\begin{document}
\begin{center}
{\fontsize{20}{23}\selectfont\bfseries Quantitative stationary crystallization\\for the $E_8$ and Leech lattices}
\end{center}
\vspace{3pt}

\section{Statement and conventions}
Let $d\in\{8,24\}$, and let $\Lambda\subset\R^d$ be the $E_8$ lattice or the Leech lattice, normalized to have covolume one. Its shortest nonzero vector length is denoted by $r_0$: thus $r_0=\sqrt2$ when $d=8$ and $r_0=2$ when $d=24$.

A configuration is a simple locally finite subset of $\R^d$. Configuration space carries the sigma algebra generated by counting points in bounded Borel sets. Let $\mu$ be a probability law on this space and write $\Xi$ for its random configuration. Stationarity means invariance of $\mu$ under every translation. Intensity one means
\[
\E_\mu\#(\Xi\cap A)=\vol(A)
\]
for every bounded Borel set $A$. Neither ergodicity nor a samplewise density limit is assumed.

Put $Q=[0,1)^d$. For $s>d$, define the ordered Riesz energy and the lattice energy by
\begin{align}
\calE_s(\mu)&=\E_\mu\sum_{x\in\Xi\cap Q}\sum_{y\in\Xi\setminus\{x\}}|x-y|^{-s},\label{eq:energy}\\
L_s&=\sum_{z\in\Lambda\setminus\{0\}}|z|^{-s}.\label{eq:lattice}
\end{align}
The second sum in \eqref{eq:energy} ranges over the entire configuration, including points outside $Q$. There is no factor of $1/2$. Local finiteness of the lattice and its polynomial point growth imply $L_s<\infty$ for $s>d$. Set
\begin{equation}
R(\delta)=\frac{\log(1/\delta)}{\log\log(1/\delta)}\qquad(0<\delta<e^{-e}).\label{eq:radius}
\end{equation}
All logarithms are natural. Balls used for observation are closed.

\begin{theorem}[Quantitative stationary crystallization]\label{thm:main}
For every $s>d$, there exist $C>0$ and $0<\delta_0<e^{-e}$, depending only on $d,s$, with the following property. If $\mu$ is stationary of intensity one, $0<\delta<\delta_0$, and
\[
\calE_s(\mu)\le L_s+\delta,
\]
then, with $R=R(\delta)$, there is a measurable event $G$ such that
\[
\mu(G)\ge1-CR^{-1/2}.
\]
For every configuration in $G$, the set $X=\Xi\cap\BR$ satisfies
\begin{gather}
\exists\ \Phi\text{ a Euclidean isometry},\quad
\exists\ Z\subset\Lambda\text{ finite}:\quad
d_H(X,\Phi Z)\le C\delta^{1/13},\label{eq:match}\\
\#X\ge(1-CR^{-1/2})\vol(\BR).\label{eq:count}
\end{gather}
\end{theorem}

The constants are chosen before $\mu$ and $\delta$. The latter is any positive upper bound on the excess, not necessarily the excess itself. Every original point of $\Xi\cap\BR$ participates in \eqref{eq:match}. The image $\Phi Z$ need not be contained in the observation ball, and no prescribed bijection is required.

For finite sets, the Hausdorff condition means that each point of either set is within the given distance of a point of the other. This also specifies the convention for empty sets. For sufficiently small $\delta$, \eqref{eq:count} is positive on $G$, so the empty-set convention does not affect the conclusion.

The proof below establishes the stronger position estimate $C\delta^\gamma$ for every fixed $0<\gamma<1/12$, with constants depending also on $\gamma$. Section 4 proves that the full event described by the two conclusions is Borel; it can therefore be used as $G$.

This manuscript presents a proposed paper proof based on two established inputs: the Gaussian auxiliary functions of Cohn, Kumar, Miller, Radchenko, and Viazovska (CKMRV), and the periodic sphere-packing stability theorem of B\"or\"oczky, Radchenko, and Ramos (BRR). Sections 2--7 supply the stationary and quantitative connections. Sections 8--9 give extensions to positive Gaussian mixtures and a classification of the equality laws.

\section{Linear programming for a stationary law}
For an even function $f$ for which the expectation is absolutely convergent, write
\[
\calP_\mu(f)=\E_\mu\sum_{x\in\Xi\cap Q}\sum_{y\in\Xi\setminus\{x\}} f(y-x).
\]
For a radial function we use the same notation with $f(|y-x|)$. Nonnegative functions may initially have infinite expectations.

\subsection{Local second moments and absolute integrability}
Let $N_Q=\#(\Xi\cap Q)$. Since the diameter of $Q$ is $\sqrt d$,
\[
N_Q(N_Q-1)\le d^{s/2}\sum_{\substack{x,y\in\Xi\cap Q\\x\ne y}}|x-y|^{-s}.
\]
The energy hypothesis and intensity one imply
\begin{equation}
M:=\E N_Q^2\le1+d^{s/2}\calE_s(\mu)<\infty.\label{eq:second}
\end{equation}
Write $N_k=\#(\Xi\cap(Q+k))$ for $k\in\Z^d$. Stationarity and Cauchy--Schwarz give $\E[N_0N_k]\le M$. Thus every Schwartz function $f$ satisfies
\begin{equation}
\E\sum_{x\in\Xi\cap Q}\sum_{y\in\Xi}|f(y-x)|
\le M\sum_{k\in\Z^d}\sup_{u\in k+[-1,1]^d}|f(u)|<\infty.\label{eq:integrability}
\end{equation}
Rapid decay makes the last sum finite. This controls signed functions and all subsequent pair-measure limits.

\subsection{Fourier positivity and the large-box limit}
We use the convention
\[
\widehat f(\xi)=\int_{\R^d}f(x)e^{-2\pi i x\cdot\xi}\,dx.
\]
Suppose $f$ is real, even, and Schwartz, with $\widehat f\ge0$. For $W_L=[0,L)^d$, introduce the finite signed measure
\[
\sigma_L=\sum_{x\in\Xi\cap W_L}\delta_x-\ind_{W_L}\,dx.
\]
Fourier inversion gives, for each configuration,
\[
\iint f(x-y)\,d\sigma_L(x)d\sigma_L(y)
=\int\widehat f(\xi)|\widehat{\sigma_L}(\xi)|^2\,d\xi\ge0.
\]
Finite local second moments justify taking expectations. The first-moment intensity formula evaluates the cross terms, leaving
\begin{equation}
\E\sum_{x,y\in\Xi\cap W_L}f(x-y)
\ge\int_{W_L}\int_{W_L}f(x-y)\,dx\,dy.\label{eq:finiteLP}
\end{equation}

For completeness, define the reduced pair measure, including the diagonal, by
\[
\int g(u)\,d\beta(u)=\E\sum_{x\in\Xi\cap Q}\sum_{y\in\Xi}g(y-x)
\]
for nonnegative Borel $g$. In coordinates $(x,u)=(x,y-x)$, the expected pair measure is $dx\otimes\beta$. Indeed, after restricting $u$ to a bounded Borel set, the resulting measure in $x$ is locally finite and translation invariant. It is a constant multiple of Lebesgue measure. Identification on $Q$ and a monotone-class argument give the product representation.

Divide \eqref{eq:finiteLP} by $|W_L|$. The left side becomes
\[
\int f(u)\frac{|W_L\cap(W_L-u)|}{|W_L|}\,d\beta(u).
\]
The weight lies in $[0,1]$ and tends to one for each fixed $u$. Equation \eqref{eq:integrability} permits dominated convergence. The same argument with Lebesgue measure treats the right side. Separating the diagonal proves the following statement.

\begin{proposition}[Stationary Schwartz bound]\label{prop:LP}
If $\mu$ is stationary of intensity one with finite local second moment, then every real even Schwartz function $f$ with $\widehat f\ge0$ satisfies
\begin{equation}
f(0)+\calP_\mu(f)\ge\widehat f(0).\label{eq:LP}
\end{equation}
\end{proposition}
No ergodic decomposition or samplewise limit enters this argument. The single diagonal term is exactly $f(0)$, consistently with the ordered normalization.

\section{A compact-parameter cost for close pairs}
\subsection{Gaussian auxiliary functions and calibration}
For $\alpha>0$, put
\[
g_\alpha(r)=e^{-\alpha r^2},\qquad
f_\alpha(x)=F(i\alpha/\pi,x),\qquad h_\alpha=g_\alpha-f_\alpha,
\]
where $F$ is the CKMRV auxiliary family. We use Lemma 6.2, Proposition 6.1, and the construction in Sections 5--6 of CKMRV. In the relevant dimension, $f_\alpha$ is real, radial, and Schwartz; its Fourier transform is nonnegative; and $h_\alpha\ge0$. For $r>0$, the zeros of $h_\alpha$ are exactly the nonzero lattice-vector lengths. The family is jointly continuous, with the compact-parameter Schwartz control supplied by that construction.

At nonzero lattice vectors, $f_\alpha=g_\alpha$ and $\widehat f_\alpha=0$. The self-duality of $\Lambda$ and Poisson summation therefore give
\begin{equation}
\widehat f_\alpha(0)-f_\alpha(0)
=\sum_{z\in\Lambda\setminus\{0\}}e^{-\alpha|z|^2}=:L_\alpha.\label{eq:calibration}
\end{equation}
Applying Proposition \ref{prop:LP} yields
\begin{equation}
\calP_\mu(g_\alpha)-L_\alpha\ge\calP_\mu(h_\alpha)\ge0
\qquad(\alpha>0).\label{eq:GaussianGap}
\end{equation}
For each fixed parameter, these expectations are finite by \eqref{eq:integrability}.

\subsection{Strict positivity down to the origin}
Put $y=\pi/\alpha$. The functional identity used in the proof of CKMRV Lemma 6.2, together with its equation (6.1), gives
\begin{align}
h_\alpha(r)&=y^{d/2}\widetilde F(iy,r),\\
h_\alpha(0)&=\frac{2\pi}{d}y^{d/2+1}
\sum_{z\in\Lambda\setminus\{0\}}|z|^2e^{-\pi y|z|^2}>0.\label{eq:origin}
\end{align}
In the Leech case, $\sqrt2$ is not a nonzero lattice-vector length, so $h_\alpha(\sqrt2)>0$ as well. Take
\[
I=[\pi/4,\pi/3].
\]
Strict positivity below $r_0$, the origin formula, and joint continuity imply
\begin{equation}
m_0:=\min_{\substack{\alpha\in I\\0\le r\le7r_0/8}}h_\alpha(r)>0.\label{eq:smallradius}
\end{equation}
This uses the regularized auxiliary function. In particular, no unregularized Laplace integral is evaluated at the origin or at the problematic small radii in dimension twenty-four.

\subsection{A quadratic bound near the shortest length}
For $7r_0/8\le r<r_0$, CKMRV equations (6.2) and (6.4)(1) give
\begin{equation}
h_\alpha(r)=4y^{d/2}\sin^2(\pi r^2/2)
\int_0^\infty\widehat K(iy,iv)e^{-\pi r^2v}\,dv.\label{eq:kernel}
\end{equation}
For $d=24$, this interval lies above $\sqrt2$, within the convergence range. For $\alpha\in I$, we have $y\in[3,4]$. The kernel is continuous and positive on the rectangle $[3,4]\times[1,2]$, which avoids the singular curves in the quotient presentation. Consequently
\[
m_K:=\min_{\substack{3\le y\le4\\1\le v\le2}}\widehat K(iy,iv)>0.
\]
Let $D=r_0^2-r^2$. In this interval $0\le D\le1$, and $r_0^2\in\{2,4\}$. Hence
\[
\sin^2(\pi r^2/2)=\sin^2(\pi D/2)\ge D^2\ge r_0^2(r_0-r)^2.
\]
Restricting the integral in \eqref{eq:kernel} to $[1,2]$, and using $r^2\le4$, proves
\begin{equation}
h_\alpha(r)\ge c_1(r_0-r)^2,\qquad
c_1=4\,3^{d/2}m_K e^{-8\pi}r_0^2>0.\label{eq:quadratic}
\end{equation}

\subsection{Integration against the Riesz weight}
Define
\[
w_s(\alpha)=\frac{\alpha^{s/2-1}}{\Gamma(s/2)},\qquad
V_s(r)=\int_I w_s(\alpha)h_\alpha(r)\,d\alpha.
\]
The Gamma-integral identity $r^{-s}=\int_0^\infty w_s(\alpha)e^{-\alpha r^2}\,d\alpha$ and Tonelli's theorem give
\[
\calE_s(\mu)=\int_0^\infty w_s(\alpha)\calP_\mu(g_\alpha)\,d\alpha,
\qquad L_s=\int_0^\infty w_s(\alpha)L_\alpha\,d\alpha.
\]
Both integrals are finite. Since \eqref{eq:GaussianGap} holds at every positive parameter, their difference is an integral of a nonnegative gap. Retaining only the cost on $I$ yields
\begin{equation}
\calP_\mu(V_s)\le\calE_s(\mu)-L_s\le\delta.\label{eq:costbudget}
\end{equation}
The inequalities outside $I$ are essential to this step. We do not treat the full positive-parameter mixture of $f_\alpha$ as a Schwartz function.

Write $w_I=\int_I w_s(\alpha)\,d\alpha>0$. Equations \eqref{eq:smallradius} and \eqref{eq:quadratic} imply, with
\[
a_0=w_I m_0,\qquad b_0=\min\{w_I m_0/r_0^2,w_I c_1\}>0,
\]
that $V_s(r)\ge a_0$ for $0<r\le r_0/2$ and $V_s(r)\ge b_0(r_0-r)^2$ for $r_0/2\le r<r_0$. We have therefore proved the precise short-distance estimate needed below.

\begin{proposition}[Short-distance cost]\label{prop:cost}
There is $\kappa=\min\{4a_0,b_0r_0^2\}>0$, depending only on $d,s$, such that
\begin{equation}
0<t\le\tfrac12,\quad0<r<r_0(1-t)
\quad\Longrightarrow\quad V_s(r)\ge\kappa t^2.\label{eq:shortcost}
\end{equation}
The same function satisfies \eqref{eq:costbudget}.
\end{proposition}

\section{Borel measurability of the local crystal event}
Let $P$ be a closed fundamental parallelepiped of $\Lambda$. Every isometric copy of the full lattice has the form $U(\Lambda+u)$ with $U\in O(d)$ and $u\in P$. The parameter space
\[
T=O(d)\times P
\]
is compact and metrizable. For a finite set $X\subset\R^d$, define
\begin{equation}
D_\Lambda(X)=\min_{(U,u)\in T}\max_{x\in X}\dist(U^Tx-u,\Lambda),\label{eq:distancefunctional}
\end{equation}
with maximum zero if $X$ is empty. The objective is continuous in $(U,u)$, so the minimum is attained.

\begin{lemma}[Finite matching and full-lattice distance]\label{lem:matching}
For every finite $X$ and $\varepsilon\ge0$,
\begin{equation}
D_\Lambda(X)\le\varepsilon
\quad\Longleftrightarrow\quad
\exists\Phi\text{ isometry},\ \exists Z\subset\Lambda\text{ finite},\quad
d_H(X,\Phi Z)\le\varepsilon.\label{eq:matchingiff}
\end{equation}
\end{lemma}
\begin{proof}
The right side puts every point of $X$ within $\varepsilon$ of a full isometric lattice, proving the forward inequality for the functional. Conversely, choose minimizing $(U,u)$ and, for each $x\in X$, a nearest lattice point $z_x$ to $U^Tx-u$. Such a point exists because $\Lambda$ is closed and discrete. Set $Z=\{z_x:x\in X\}$ and $\Phi(z)=U(z+u)$. Each point of $X$ is close to its selected image, and every image point was selected by some point of $X$. Repeated choices cause no problem. If $X$ is empty, take $Z$ empty.
\end{proof}

Fix a countable dense subset $\{(U_j,u_j):j\ge1\}$ of $T$. For $X=\Xi\cap\BR$,
\[
D_\Lambda(X)=\inf_j M_j(\Xi),\qquad
M_j(\Xi)=\max_{x\in\Xi\cap\BR}\dist(U_j^Tx-u_j,\Lambda).
\]
For $v\ge0$,
\[
\{M_j>v\}=\left\{\Xi\!\left(\BR\cap
\{x:\dist(U_j^Tx-u_j,\Lambda)>v\}\right)\ge1\right\}.
\]
The set inside the count is bounded and Borel, so $M_j$ is measurable. Thus $D_\Lambda(\Xi\cap\BR)$ is measurable. The point-count condition is measurable as well.

\begin{proposition}\label{prop:Borel}
For fixed $R,\varepsilon$ and a fixed point-count threshold, the event consisting of finite-subset matching within $\varepsilon$ and the point-count lower bound is Borel. The corresponding event with a variable observation center is jointly Borel.
\end{proposition}
The last assertion follows by applying the same argument to the jointly measurable translation map and parameterized counts. No continuity as a point crosses the boundary of the closed ball is needed. No measurable selection of $\Phi$ or $Z$, and no projection of an uncountable family of finite witnesses, is used.

\section{From periodic stability to stationary stability}
We first record the periodic input in the form required here. Theorems 1.5 and 1.6 of BRR apply to minimum-distance-$r_0$ periodic packings of center density at least $1-\eta$. Put
\[
R_\eta=R(\eta),\qquad b_\eta=B_dR_\eta^{-1/2},\qquad
\varepsilon_\eta=E_d e^{c_dR_\eta}\eta^{1/4},
\]
where the positive dimensional constants incorporate those of BRR. For a centered closed ball $K=\overline B_r$, the admissibility conditions are
\begin{equation}
R_\eta\le r,\qquad 2r\le M_dR_\eta,
\quad M_8=2^{20},\quad M_{24}=2^{140}.\label{eq:admissible}
\end{equation}
For sufficiently small $\eta$, a uniform observation center modulo the period lattice then has, with probability at least $1-b_\eta$, both a finite-subset match of all points in the window with error $\varepsilon_\eta$ and at least $(1-b_\eta)\vol(K)$ points. Theorems 1.1 and 1.2 recorded in BRR also provide the periodic center-density bound one in these normalizations.

\begin{lemma}[Stationary transfer]\label{lem:transfer}
Let $Y$ be a stationary point process of finite intensity $\rho$, whose distinct points have distance at least $r_0$. Then $\rho\le1$. For admissible $\eta,r$, the probability that the two periodic-stability conclusions fail for $Y\cap\overline B_r$ is at most
\begin{equation}
b_\eta+\frac{1-\rho}{\eta}.\label{eq:transfer}
\end{equation}
\end{lemma}
\begin{proof}
Set $h=r_0$. For $L>2h$ let
\[
W_L=[0,L)^d,\qquad W_L^-=[h,L-h)^d,\qquad
P_L(Y)=(Y\cap W_L^-)+L\Z^d.
\]
Points from the same cell retain the original separation. Points from different cells have a gap of at least $2h$ in some coordinate. Thus $P_L(Y)$ is a periodic packing of minimum distance at least $r_0$. Its center density satisfies
\[
D_L(Y)=\frac{\#(Y\cap W_L^-)}{L^d}\le1.
\]
The intensity identity gives
\begin{equation}
\E D_L=\rho(1-2h/L)^d.\label{eq:meandensity}
\end{equation}
Letting $L$ increase first shows $\rho\le1$. Since $1-D_L\ge0$, Markov's inequality implies
\begin{equation}
\Prob(D_L<1-\eta)
\le\frac{1-\rho(1-2h/L)^d}{\eta}.\label{eq:markovdensity}
\end{equation}

Let $U_L$ be uniform on $W_L$, independent of $Y$. For every sample with $D_L\ge1-\eta$, apply the periodic theorem at $U_L$. Proposition \ref{prop:Borel} justifies integration on the product space. Periodization is measurable: its count in a bounded Borel set is a countable sum of counts in translated Borel sets inside $W_L^-$. Thus the probability of failure for the periodized process is at most the right side of \eqref{eq:markovdensity} plus $b_\eta$.

If $U_L$ is at distance more than $h+r$ from the boundary of $W_L$, then $Y$ and $P_L(Y)$ agree on $U_L+\overline B_r$. The fraction of excluded centers is at most
\[
1-\left(1-\frac{2(h+r)}{L}\right)^d=O_d((h+r)/L)
\]
for sufficiently large $L$. Stationarity says that $Y-U_L$ has the same law as $Y$. Therefore, for the fixed original process and window,
\[
\Prob(\mathrm{failure})\le b_\eta+
\frac{1-\rho(1-2h/L)^d}{\eta}+O_d((h+r)/L).
\]
The left side is independent of $L$. Letting $L\to\infty$ proves \eqref{eq:transfer}.
\end{proof}
Only expectations and explicit boundary errors have been taken to a limit. There is no assertion that the density of an individual sample converges. Equivalently, one can choose a finite common $L$ that makes both additional errors smaller than any prescribed positive tolerance.

\section{Thinning, dilation, and restoration of the original points}
Fix $0<t\le1/2$. Delete every point of $\Xi$ that has another point at distance strictly less than $r_0(1-t)$, and call the remaining process $Y_t$. This decision is made in the whole configuration, so neighbors outside the observation window are included. The deletion map is measurable by parameterized counting, and it commutes with translations. Consequently $Y_t$ and the deleted process are stationary.

Let $q_t$ be the intensity of the deleted process. Every deleted point has at least one offending neighbor. Proposition \ref{prop:cost} therefore gives
\begin{align}
q_t&\le\E\sum_{x\in\Xi\cap Q}\sum_{y\ne x}
\ind_{\{|x-y|<r_0(1-t)\}}\notag\\
&\le\frac{\calP_\mu(V_s)}{\kappa t^2}
\le\frac{\delta}{\kappa t^2}.\label{eq:deleted}
\end{align}
The retained intensity is $1-q_t$, and its minimum distance is at least $r_0(1-t)$. Ordered pairs are used throughout; no factor of $1/2$ appears.

Dilate by
\[
\lambda=(1-t)^{-1},\qquad Y'_t=\lambda Y_t.
\]
This is a stationary packing of minimum distance $r_0$. Its intensity is
\begin{equation}
\rho_t=(1-q_t)(1-t)^d,\qquad
1-\rho_t\le q_t+dt.\label{eq:dilatedintensity}
\end{equation}
Apply Lemma \ref{lem:transfer} on the ball of radius $\lambda R$, assuming the parameters meet \eqref{eq:admissible}. Outside its failure event, and for $b_\eta<1$,
\begin{align}
\#(Y_t\cap\BR)&=\#(Y'_t\cap\overline B_{\lambda R})\notag\\
&\ge(1-b_\eta)\lambda^d\vol(\BR)
\ge(1-b_\eta)\vol(\BR).\label{eq:restoredcount}
\end{align}

The position estimate must retain an actual Euclidean isometry of the unscaled lattice. Put $X_t=Y_t\cap\BR$. Every point $x\in X_t$ is within $(\lambda-1)R\le2tR$ of $\lambda x$. Thus
\[
d_H(X_t,\lambda X_t)\le2tR.
\]
If $\Phi,Z$ are supplied by the stability theorem for $\lambda X_t$, the same witnesses give
\begin{equation}
d_H(X_t,\Phi Z)\le\varepsilon_\eta+2tR.\label{eq:restorederror}
\end{equation}
The map $\Phi$ itself has not been rescaled.

Finally, the intensity formula for the deleted process and Markov's inequality show that
\begin{equation}
\Prob\big((\Xi\setminus Y_t)\cap\BR\ne\varnothing\big)
\le q_t\vol(\BR).\label{eq:nodeletions}
\end{equation}
On the complementary event, $X_t=\Xi\cap\BR$. Taking a union bound, without any independence assumption, proves the combined estimate.

\begin{proposition}[Combined quantitative bound]\label{prop:combined}
Whenever the above parameters are admissible, the conclusions
\[
d_H(\Xi\cap\BR,\Phi Z)\le\varepsilon_\eta+2tR,
\qquad
\#(\Xi\cap\BR)\ge(1-b_\eta)\vol(\BR)
\]
hold for suitable $\Phi,Z$ outside an event of probability at most
\begin{equation}
b_\eta+\frac{dt+\delta/(\kappa t^2)}{\eta}
+\frac{\delta}{\kappa t^2}\vol(\BR).\label{eq:failurebound}
\end{equation}
\end{proposition}
This is the step that restores all original points in the observation ball; the conclusion is not restricted to a selected subset of them.

\section{Parameter choice and completion of the proof}
Fix $0<\gamma<1/12$ and set
\begin{equation}
t=\delta^{1/3},\qquad
\beta=\frac{1+12\gamma}{6}\in(1/6,1/3),\qquad
\eta=\delta^\beta.\label{eq:parameters}
\end{equation}
Write $\ell=\log(1/\delta)$. Then
\[
R=\frac{\ell}{\log\ell},\qquad R_\eta=\frac{\beta\ell}{\log(\beta\ell)}.
\]
For sufficiently small $\delta$, we have $\beta\ell>e$ and
\begin{equation}
\beta R\le R_\eta\le R.\label{eq:radiuscomparison}
\end{equation}
Indeed, the lower bound uses $\log(\beta\ell)\le\log\ell$, and the upper bound uses monotonicity of $x/\log x$ for $x>e$. Since $\lambda\le2$,
\[
R_\eta\le\lambda R,\qquad
\frac{2\lambda R}{R_\eta}\le\frac4\beta<24.
\]
Both dimensional diameter conditions in \eqref{eq:admissible} follow.

For the position estimate,
\[
\varepsilon_\eta=E_d\exp(c_dR_\eta)\delta^{\beta/4}
=\delta^{\beta/4-o(1)}.
\]
The positive margin is
\begin{equation}
\frac\beta4-\gamma=\frac{1-12\gamma}{24}>0.\label{eq:positionmargin}
\end{equation}
Thus $\varepsilon_\eta\le C\delta^\gamma$ after reducing the threshold. Also $2tR=2\delta^{1/3}R=o(\delta^\gamma)$.

For the probability bound, \eqref{eq:parameters} gives $\delta/t^2=\delta^{1/3}$ and
\begin{gather}
\frac{t}{\eta}=\delta^{1/3-\beta}
=\delta^{(1-12\gamma)/6}=o(R^{-1/2}),\label{eq:probmargin}\\
\delta^{1/3}R^d=o(R^{-1/2}),\qquad
R_\eta^{-1/2}\le\beta^{-1/2}R^{-1/2}.\label{eq:probdecay}
\end{gather}
To see these comparisons directly, $R=\ell/\log\ell$ grows subexponentially in $\ell$, whereas each fixed positive power of $\delta$ is $e^{-a\ell}$. Substitution in \eqref{eq:failurebound} proves the failure bound $CR^{-1/2}$. Equation \eqref{eq:restoredcount} gives the same order of point-count loss.

Choose $\delta_0$ to satisfy the finitely many smallness conditions, including the BRR threshold, $t\le1/2$, $b_\eta<1$, and the estimates just proved. Increase $C$ to dominate the finitely many constants. Both choices depend only on $d,s,\gamma$, and precede the choice of $\mu$. Proposition \ref{prop:Borel} makes the full success event Borel. This proves the stronger statement below, and Theorem \ref{thm:main} follows by specialization.

\begin{theorem}[Range of exponents]\label{thm:gamma}
For every fixed $0<\gamma<1/12$, Theorem \ref{thm:main} holds with $C\delta^\gamma$ in place of $C\delta^{1/13}$, with constants depending on $d,s,\gamma$. The full geometric success event is Borel.
\end{theorem}
For $\gamma=1/13$, the explicit choice is $\beta=25/78$, and
\[
\frac\beta4-\frac1{13}=\frac1{312},\qquad
\frac13-\beta=\frac1{78}.
\]
The endpoint $\gamma=1/12$ is not obtained by this argument: both strict exponent margins disappear. No assertion about the truth or falsehood of the endpoint is made.

\section{Nonzero positive Gaussian mixtures}
Let $\nu$ be a nonzero positive Borel measure on $(0,\infty)$, and put
\begin{equation}
p(r)=\int_{(0,\infty)}e^{-\alpha r^2}\,d\nu(\alpha),\qquad
L_p=\sum_{z\in\Lambda\setminus\{0\}}p(|z|)<\infty.\label{eq:mixture}
\end{equation}
Define $\calE_p(\mu)$ by replacing $|x-y|^{-s}$ in \eqref{eq:energy} with $p(|x-y|)$. For an ordinary finite-valued potential one may additionally impose $p(r)<\infty$ for all $r>0$. The argument also permits extended values: the finite-energy hypothesis then excludes pairs at distances of infinite cost.

\begin{theorem}[Positive mixtures]\label{thm:mixture}
For every potential \eqref{eq:mixture} and every fixed $0<\gamma<1/12$, the conclusion of Theorem \ref{thm:gamma} holds under
\[
\calE_p(\mu)\le L_p+\delta,
\]
with constants depending on $d,p,\gamma$. The measure $\nu$ need not have a density; in particular, a single Gaussian is allowed.
\end{theorem}
\begin{proof}
Finiteness of $L_p$ implies $p(r_0)<\infty$. For $0<a\le b<\infty$,
\[
\nu([a,b])e^{-br_0^2}\le p(r_0)<\infty.
\]
Thus $\nu$ is locally finite and sigma finite. Since it is nonzero, some compact interval $I=[a,b]\subset(0,\infty)$ has finite positive mass. For $0<r\le D<\infty$,
\[
p(r)\ge\nu(I)e^{-bD^2}>0.
\]
With $D=\sqrt d$, this lower bound and finite energy yield $\E N_Q^2<\infty$, by the same pair-count estimate as in Section 2. Proposition \ref{prop:LP} and hence \eqref{eq:GaussianGap} apply to every $\alpha>0$.

Define the compact-parameter slack
\[
V_p(r)=\int_I h_\alpha(r)\,d\nu(\alpha).
\]
Integrate \eqref{eq:GaussianGap} over all positive parameters and apply Tonelli to the potential and lattice sums. Both total energies are finite, and all Gaussian gaps are nonnegative. Therefore
\begin{equation}
\calP_\mu(V_p)\le\calE_p(\mu)-L_p\le\delta.\label{eq:mixturebudget}
\end{equation}
The strict positivity and continuity used in Section 3.2 hold on the chosen compact interval, giving the same positive small-radius lower bound.

For the quadratic estimate near $r_0$, the parameter $y=\pi/\alpha$ lies in $[\pi/b,\pi/a]$. In \eqref{eq:kernel}, restrict the integration variable to $[V,V+1]$, where
\[
V>\max\{1,\pi/a,b/\pi\}+1.
\]
This rectangle avoids both $y=v$ and $y=1/v$. The CKMRV kernel is continuous and strictly positive there, so it has a positive minimum. The exponential factor is bounded below by $e^{-\pi r_0^2(V+1)}$. The same sine-squared estimate as before gives a uniform quadratic lower bound. Integration over the finite positive mass of $I$ yields
\[
0<t\le\tfrac12,\quad0<r<r_0(1-t)
\quad\Longrightarrow\quad V_p(r)\ge\kappa_p t^2
\]
for some $\kappa_p>0$. Sections 4--7 now apply without change.
\end{proof}

\section{Classification of the zero-excess laws}
In this section the potential is either $r^{-s}$ with $s>d$, or a nonzero positive Gaussian mixture satisfying \eqref{eq:mixture}. Write $\calE_p,L_p$ for its energy and lattice value in either case. The orientation space is $O(d)/\Aut(\Lambda)$, where $\Aut(\Lambda)$ is the orthogonal automorphism group of the lattice.

\begin{theorem}[Equality laws]\label{thm:equality}
A stationary simple locally finite point process of intensity one has $\calE_p(\mu)=L_p$ if and only if its law is a mixture of isometric lattice laws with Haar-uniform translation phase. More precisely,
\begin{equation}
\mu=\int_{O(d)/\Aut(\Lambda)}
\Law\bigl(U(\Lambda+T)\bigr)\,d\sigma([U]),
\qquad T\sim\Haar(\R^d/\Lambda),\label{eq:equalitylaw}
\end{equation}
for some probability measure $\sigma$ on the orientation space.
\end{theorem}

\subsection{An almost-sure containing lattice}
Assume $\calE_p(\mu)=L_p$. Since every positive $\delta$ is an upper bound on the excess, apply the quantitative theorem with $\gamma=1/13$ and
\[
\delta_n=e^{-n^8},\qquad R_n=\frac{n^8}{8\log n}\quad(n\ge2).
\]
The thresholds hold for all sufficiently large $n$, and $\sum_nR_n^{-1/2}<\infty$. Borel--Cantelli implies that almost every configuration satisfies, for all sufficiently large $n$,
\begin{gather}
\Xi\cap\overline B_{R_n}\text{ lies within }C\delta_n^{1/13}
\text{ of }\Gamma_n=U_n(\Lambda+u_n),\label{eq:containapprox}\\
\#(\Xi\cap\overline B_{R_n})
\ge(1-CR_n^{-1/2})\vol(\overline B_{R_n}).\label{eq:ascount}
\end{gather}
Normalize $U_n\in O(d)$ and $u_n\in P$. By compactness, a subsequence converges to $(U,u)$. Every fixed point $x\in\Xi$ eventually lies in each of the expanding balls. Distance to the closed lattice is continuous, so \eqref{eq:containapprox} implies
\begin{equation}
\Xi\subset\Gamma:=U(\Lambda+u).\label{eq:contain}
\end{equation}

\subsection{Uniqueness and measurability of the containing lattice}
The intersection of two affine lattices, if nonempty, is a coset of the intersection of their vector lattices. If the latter has rank less than $d$, the intersection has density zero. If it has rank $d$, its index in either covolume-one lattice is a positive integer. Unless the affine lattices coincide, this index is at least two, and the intersection has density at most $1/2$.

Equation \eqref{eq:ascount} rules out containment of $\Xi$ in the intersection of two distinct covolume-one affine lattices. Hence $\Gamma$ in \eqref{eq:contain} is unique for almost every configuration.

Here is a direct measurability justification. Let $\mathscr L$ be the space of all isometric images of $\Lambda$, regarded as point-counting measures with the vague topology. It is a compact metrizable space: it is the image of $O(d)\times P$, and for a compactly supported continuous test function only uniformly finitely many lattice points contribute to its sum. Thus that map is continuous.

The relation $\Xi\subset\Gamma$ is the measure inequality $\Xi\le\Gamma$, a closed relation in the vague topology. Projection over a compact subset of $\mathscr L$ preserves closedness. Nonuniqueness is Borel as well: it is the countable union, over positive integers $m$, of the closed projections asserting two containing lattices at distance at least $1/m$. The usual counting sigma algebra for locally finite configurations agrees with the Borel sigma algebra of this topology.

On the Borel set $A$ of configurations with a unique containing lattice, for every closed $C\subset\mathscr L$,
\[
\{\Xi\in A:\Gamma(\Xi)\in C\}
=A\cap\{\Xi:\exists\Gamma\in C,\ \Xi\subset\Gamma\}.
\]
This proves that $\Xi\mapsto\Gamma(\Xi)$ is Borel on $A$. The set $A$ is translation invariant, and uniqueness gives
\[
\Gamma(\Xi+a)=\Gamma(\Xi)+a.
\]
Its pushforward under $\mu$, which gives $A$ full measure, is therefore a stationary lattice process.

\subsection{Expected intensity removes all missing points}
A bounded fundamental domain of volume one yields the uniform lattice count
\begin{equation}
\#(\Gamma\cap\overline B_R)
=\vol(\overline B_R)+O_\Lambda(R^{d-1})\label{eq:latticecount}
\end{equation}
over all isometric images $\Gamma$. To obtain it, compare the volume-one cells associated with the counted lattice points to the balls enlarged and reduced by the diameter of the fundamental domain. Stationarity and \eqref{eq:latticecount} show that the intensity of the random containing lattice is one.

Both $\Xi$ and $\Gamma$ thus have intensity one, and $\Xi\subset\Gamma$. For each $k\in\Z^d$,
\[
\E\#\bigl((\Gamma\setminus\Xi)\cap(Q+k)\bigr)=0.
\]
The count is a nonnegative integer, so it vanishes almost surely. Taking the intersection of these probability-one events over countably many cubes proves $\Xi=\Gamma$ almost surely. This excludes even a set of missing points of density zero.

\subsection{Haar phase and arbitrary orientation mixture}
Let $\sigma$ be the distribution of the orientation of $\Gamma$. For every continuous function $H$ on the compact space $\mathscr L$, stationarity gives
\[
\E H(\Gamma)=\E\left[\frac1{|W_T|}\int_{W_T}H(\Gamma+a)\,da\right].
\]
For a fixed lattice, tiling by a bounded fundamental domain and controlling the boundary cells show that the inner average tends to the Haar average over its phase torus. It depends only on the orientation. Dominated convergence therefore proves \eqref{eq:equalitylaw}; continuous functions determine probability measures on the compact space.

Conversely, Haar phase makes each fixed-orientation lattice law stationary of intensity one. At every lattice point its off-diagonal potential sum is $L_p$, so its ordered expected energy is $L_p$. Mixing orientations preserves these properties. This completes the proof of Theorem \ref{thm:equality}. No rotational invariance was assumed; the orientation measure $\sigma$ is arbitrary.
\section{Conclusion}
The argument gives quantitative local crystallization for stationary intensity-one point processes in dimensions eight and twenty-four. An ordered Riesz-energy excess of at most $\delta$ yields position error $C\delta^{1/13}$ on the closed ball of radius $R=\log(1/\delta)/\log\log(1/\delta)$, with probability and point-count losses at most $CR^{-1/2}$. The stationary Fourier bound and compact-parameter short-distance cost control close pairs; expected density transfers periodic stability; and the no-deletion estimate restores every original point in the observation ball. The lattice witness remains a Euclidean isometric image of a finite subset of the covolume-one $E_8$ or Leech lattice.

The same estimates give every fixed position exponent below $1/12$ and apply to nonzero positive Gaussian mixtures of finite lattice energy. At zero excess, compactness produces a unique containing lattice, intensity one removes every missing point, and stationarity forces Haar-uniform translation phase. Thus the equality laws are precisely the orientation mixtures described in Theorem \ref{thm:equality}.
\end{document}
